\chapter{Complexity — Wrappers}
%%% generated by the blueprint pipeline from TCSlib.Complexity.TuringMachine.Build.Wrappers
%%%%%%%%%%%%%%%%%%%%%%%%%%%%%%%%%%%%%%%%%%%%%%%%%%%%%%%%%%%%%%%%%%%%%%%%%%%%%%%
\section{Overview}

This module provides the output-isolation wrappers of the machine-construction library
for multi-tape Turing machines over the Boolean alphabet. It contains the capture
transform, under which a host machine simulates a source machine in lockstep while
recording the source's output on one extra work tape and emitting nothing itself; its
forwarding dual, under which the source's output is appended to the host's own output;
the halt-redirect transformation, which runs a machine silently and halts exactly when
the last bit it emitted is a designated bit, entering an idle loop otherwise; and a timed
conditional, which computes $x \mapsto f_1(x)$ if $p(x)$ and $x \mapsto f_2(x)$ otherwise
within a constant multiple of the decider's running time plus the slower branch's
running time.

%%%%%%%%%%%%%%%%%%%%%%%%%%%%%%%%%%%%%%%%%%%%%%%%%%%%%%%%%%%%%%%%%%%%%%%%%%%%%%%
\section{Declarations}

\begin{definition}[Capture transform of a machine action]
\lean{Turing.captureAction}
\label{Turing.captureAction}
\leanok
\uses{Turing.Action}
Let $S$ and $H$ be sets of control states, $\iota : S \to H$ a map and $r \in H$ a
designated return state. Given an action $a$ of a Boolean Turing machine with $k$ work
tapes and state set $S$, its capture transform is the action of a machine with $k+1$
work tapes and state set $H$ that moves the input head as $a$ does and performs $a$'s
write and head move on each of the first $k$ work tapes; on the last work tape it writes
the bit $b$ and moves right if $a$ emits $b$, and neither writes nor moves if $a$ emits
nothing; it emits nothing on the output tape; and its successor state is $\iota(q')$ if
$a$ has successor state $q'$, and $r$ if $a$ halts.
\end{definition}

\begin{definition}[Capture image of a configuration]
\lean{Turing.captureCfg}
\label{Turing.captureCfg}
\leanok
\uses{Turing.Cfg, Turing.FinTM.bufferTape}
Let $S$ and $H$ be sets of control states, $\iota : S \to H$ a map, $r \in H$ a return
state, and $p$, $o$ Boolean words. For a configuration $c$ of a Boolean Turing machine
with $k$ work tapes and state set $S$ on an input $x$, with output word $u$, its capture
image is the configuration on $x$ of a machine with $k+1$ work tapes and state set $H$
whose control state is $\iota(q)$ if $c$ is in state $q$ and $r$ if $c$ is halted; whose
input head and first $k$ work tapes (contents and head positions) are those of $c$; whose
last work tape holds the word $pu$ in cells $0, \dots, |pu|-1$, blank elsewhere, with its
head at cell $|pu|$; and whose output tape holds $o$.
\end{definition}

\begin{lemma}[Capture transform commutes with applying an action]
\lean{Turing.capture_apply}
\label{Turing.capture_apply}
\leanok
\uses{Turing.Action, Turing.Action.apply, Turing.Cfg, Turing.FinTM.bufferTape_append, Turing.captureAction, Turing.captureCfg}
\difficulty{4}
Let $S$ and $H$ be sets of states, $\iota : S \to H$, $r \in H$, and $p$, $o$ Boolean
words. For every action $a$ and every configuration $c$ of a Boolean Turing machine with
$k$ work tapes and state set $S$ on an input $x$, applying the capture transform of $a$
to the capture image of $c$ gives the capture image of the configuration obtained by
applying $a$ to $c$. Here the capture transform of $a$ performs $a$ on the input and the
first $k$ of $k+1$ work tapes, writes $a$'s emitted bit (if any) on the last tape and
moves that head right, emits nothing, and maps the successor through $\iota$ with a halt
replaced by $r$; and the capture image of $c$ has state $\iota(q)$ (or $r$ if $c$ is
halted), the input head and work tapes of $c$ on the first $k$ tapes, the word $p$
followed by the output of $c$ stored from cell $0$ of the last tape with that head just
past its end, and output $o$.
\end{lemma}

\begin{theorem}[Capture contract: lockstep run with captured output]
\lean{Turing.capture_run}
\label{Turing.capture_run}
\leanok
\uses{Turing.Cfg, Turing.Cfg.Halted, Turing.Cfg.inputSymbol, Turing.Cfg.workTapeSymbols, Turing.MultiTapeTM, Turing.MultiTapeTM.runFrom, Turing.MultiTapeTM.runFrom_succ_eq_step', Turing.MultiTapeTM.step, Turing.captureAction, Turing.captureCfg, Turing.capture_apply}
\difficulty{5}
Let $M$ be a Turing machine over the Boolean alphabet with $k$ work tapes and state set
$S$, let $N$ be one with $k+1$ work tapes and state set $H$, and let $\iota : S \to H$
and $r \in H$. Suppose that for every $s \in S$, every (possibly blank) input symbol
$\sigma$ and every tuple $(w_0, \dots, w_k)$ of work-tape symbols, the transition of $N$
at $(\iota(s), \sigma, w_0, \dots, w_k)$ is the capture transform of the transition of
$M$ at $(s, \sigma, w_0, \dots, w_{k-1})$: $M$'s action on the input and first $k$ tapes,
its emitted bit (if any) written on the last tape followed by a right move, no output,
and a halt replaced by entering $r$. Let $p$, $o$ be Boolean words, $c_0$ a
configuration of $M$ on an input $x$, and $t \in \mathbb{N}$ such that $M$ run from $c_0$
is not halted at any time $t' < t$; then running $N$ for $t$ steps from the capture image
of $c_0$ yields the capture image of the configuration $M$ reaches after $t$ steps from
$c_0$, where the capture image of a configuration $c$ has state $\iota(q)$ (or $r$ if $c$
is halted), the input head and work tapes of $c$ on the first $k$ tapes, the word $p$
followed by the output of $c$ stored from cell $0$ of the last tape with that head just
past its end, and output $o$.
\end{theorem}

\begin{definition}[Forwarding transform of a machine action]
\lean{Turing.emitAction}
\label{Turing.emitAction}
\leanok
\uses{Turing.Action}
Let $S$ and $H$ be sets of control states, $\iota : S \to H$ a map and $r \in H$ a
designated return state. Given an action $a$ of a Boolean Turing machine with $k$ work
tapes and state set $S$, its forwarding transform is the action of a machine with the
same $k$ work tapes and state set $H$ that has the same input-head move, the same writes
and head moves on every work tape, and the same emitted bit (if any) as $a$, and whose
successor state is $\iota(q')$ if $a$ has successor state $q'$, and $r$ if $a$ halts.
\end{definition}

\begin{definition}[Forwarding image of a configuration]
\lean{Turing.emitCfg}
\label{Turing.emitCfg}
\leanok
\uses{Turing.Cfg}
Let $S$ and $H$ be sets of control states, $\iota : S \to H$ a map, $r \in H$ a return
state, and $p$ a Boolean word. For a configuration $c$ of a Boolean Turing machine with
$k$ work tapes and state set $S$ on an input $x$, its forwarding image is the
configuration on $x$ of a machine with $k$ work tapes and state set $H$ whose control
state is $\iota(q)$ if $c$ is in state $q$ and $r$ if $c$ is halted, whose input-head
position, work-tape contents and work-head positions are those of $c$, and whose output
tape holds $p$ followed by the output of $c$.
\end{definition}

\begin{lemma}[Forwarding transform commutes with applying an action]
\lean{Turing.emit_apply}
\label{Turing.emit_apply}
\leanok
\uses{Turing.Action, Turing.Action.apply, Turing.Cfg, Turing.emitAction, Turing.emitCfg}
\difficulty{2}
Let $S$ and $H$ be sets of states, $\iota : S \to H$, $r \in H$, and $p$ a Boolean word.
For every action $a$ and every configuration $c$ of a Boolean Turing machine with $k$
work tapes and state set $S$ on an input $x$, applying the forwarding transform of $a$
(same head moves, writes and emission as $a$, successor mapped through $\iota$, a halt
replaced by $r$) to the forwarding image of $c$ (state mapped through $\iota$, or $r$ if
$c$ is halted, the same tapes and heads, output $p$ followed by the output of $c$) gives
the forwarding image of the configuration obtained by applying $a$ to $c$.
\end{lemma}

\begin{theorem}[Forwarding contract: lockstep run with forwarded output]
\lean{Turing.emit_run}
\label{Turing.emit_run}
\leanok
\uses{Turing.Cfg, Turing.Cfg.Halted, Turing.Cfg.inputSymbol, Turing.Cfg.workTapeSymbols, Turing.MultiTapeTM, Turing.MultiTapeTM.runFrom, Turing.MultiTapeTM.runFrom_succ_eq_step', Turing.MultiTapeTM.step, Turing.emitAction, Turing.emitCfg, Turing.emit_apply}
\difficulty{5}
Let $M$ and $N$ be Turing machines over the Boolean alphabet with $k$ work tapes each and
state sets $S$ and $H$, and let $\iota : S \to H$ and $r \in H$. Suppose that for every
$s \in S$, every (possibly blank) input symbol $\sigma$ and every tuple $w$ of $k$
work-tape symbols, the transition of $N$ at $(\iota(s), \sigma, w)$ is the forwarding
transform of the transition of $M$ at $(s, \sigma, w)$: the same head moves, writes and
emission, with the successor mapped through $\iota$ and a halt replaced by entering $r$.
Let $p$ be a Boolean word, $c_0$ a configuration of $M$ on an input $x$, and
$t \in \mathbb{N}$ such that $M$ run from $c_0$ is not halted at any time $t' < t$; then
running $N$ for $t$ steps from the forwarding image of $c_0$ yields the forwarding image
of the configuration $M$ reaches after $t$ steps from $c_0$, where the forwarding image of
a configuration $c$ has state $\iota(q)$ (or $r$ if $c$ is halted), the input head, work
tapes and work heads of $c$, and output $p$ followed by the output of $c$.
\end{theorem}

\begin{definition}[Halt-redirected machine]
\lean{Turing.FinTM.redirectTM}
\label{Turing.FinTM.redirectTM}
\leanok
\uses{Turing.FinTM}
Let $M$ be a Turing machine over the Boolean alphabet with finitely many states, state
set $Q$, initial state $q_0$ and $k$ work tapes, and let $b$ be a bit. The halt-redirected
machine of $M$ and $b$ has $k$ work tapes, control states the pairs $(s, \rho)$ with
$s \in Q$ and $\rho$ either a bit or empty, together with one loop state $\ell$, and
initial state $(q_0, \text{empty})$. In state $(s, \rho)$ it performs $M$'s input move
and work-tape writes and moves at $s$ but emits nothing, updates the register to $\rho'$,
the bit $M$ emits if there is one and $\rho$ otherwise, and enters $(s', \rho')$ if $M$
moves to state $s'$, while if $M$ halts it halts when $\rho' = b$ and enters $\ell$
otherwise; in state $\ell$ it moves no head, writes and emits nothing, and stays in
$\ell$.
\end{definition}

\begin{definition}[Redirected control state]
\lean{Turing.FinTM.redirectState}
\label{Turing.FinTM.redirectState}
\leanok
Let $S$ be a set of states, $b$ a bit, $q$ either a state in $S$ or the halting marker,
and $\rho$ either a bit or empty. The redirected control state, an optional element of
$(S \times \{\text{empty}, \mathrm{false}, \mathrm{true}\}) \sqcup \{\ell\}$, is the
simulation state $(s, \rho)$ if $q = s \in S$; if $q$ is the halting marker, it is the
halting marker when $\rho = b$ (so never when $\rho$ is empty) and the loop state $\ell$
otherwise.
\end{definition}

\begin{definition}[Redirected action with a last-emission register]
\lean{Turing.FinTM.redirectAction}
\label{Turing.FinTM.redirectAction}
\leanok
\uses{Turing.Action, Turing.FinTM.redirectState}
Let $S$ be a set of states, $b$ a bit, $a$ an action of a Boolean Turing machine with $k$
work tapes and state set $S$, and $\rho$ either a bit or empty. Let $\rho'$ be the bit
$a$ emits if there is one, and $\rho$ otherwise. The redirected action is the action on
$k$ work tapes with state set $(S \times \{\text{empty}, \mathrm{false},
\mathrm{true}\}) \sqcup \{\ell\}$ that has the same input move and work-tape writes and
moves as $a$, emits nothing, and has successor $(s', \rho')$ if $a$ has successor $s'$,
while if $a$ halts it halts when $\rho' = b$ and enters the loop state $\ell$ otherwise.
\end{definition}

\begin{definition}[Redirected image of a configuration]
\lean{Turing.FinTM.redirectCfg}
\label{Turing.FinTM.redirectCfg}
\leanok
\uses{Turing.Cfg, Turing.FinTM, Turing.FinTM.redirectState, Turing.FinTM.redirectTM}
Let $M$ be a finite-state Turing machine over the Boolean alphabet and $b$ a bit, and
consider the halt-redirected machine of $M$ and $b$, which simulates $M$ silently while
recording the last bit $M$ has emitted, halts when $M$ halts with that bit equal to $b$,
and otherwise enters an idle loop state $\ell$. For a configuration $c$ of $M$ on an input
$x$, let $\rho$ be the last bit of the output of $c$ (empty if that output is empty); the
redirected image of $c$ is the configuration of the halt-redirected machine on $x$ with
the input head, work tapes and work heads of $c$, empty output, and control state
$(q, \rho)$ if $c$ is in state $q$, while if $c$ is halted it is halted when $\rho = b$
and in $\ell$ otherwise.
\end{definition}

\begin{lemma}[The redirect loop state is a fixed point]
\lean{Turing.FinTM.redirect_loop}
\label{Turing.FinTM.redirect_loop}
\leanok
\uses{Turing.Action.apply, Turing.Cfg, Turing.FinTM, Turing.FinTM.redirectTM, Turing.MultiTapeTM.runFrom, Turing.MultiTapeTM.runFrom_succ_eq_step', Turing.MultiTapeTM.step}
\difficulty{3}
Let $M$ be a finite-state Turing machine over the Boolean alphabet and $b$ a bit, and
consider the halt-redirected machine of $M$ and $b$, which simulates $M$ silently while
recording the last bit $M$ has emitted, halts when $M$ halts with that bit equal to $b$,
and otherwise enters a loop state $\ell$. For every configuration $c$ of the
halt-redirected machine on an input $x$ whose control state is $\ell$, and every
$t \in \mathbb{N}$, running the halt-redirected machine for $t$ steps from $c$ returns
exactly $c$.
\end{lemma}

\begin{lemma}[Redirection commutes with applying an action]
\lean{Turing.FinTM.redirect_apply}
\label{Turing.FinTM.redirect_apply}
\leanok
\uses{Turing.Action, Turing.Action.apply, Turing.Cfg, Turing.FinTM, Turing.FinTM.redirectAction, Turing.FinTM.redirectCfg}
\difficulty{3}
Let $M$ be a finite-state Turing machine over the Boolean alphabet with $k$ work tapes,
$b$ a bit, $c$ a configuration of $M$ on an input $x$, $a$ an action of $M$, and $\rho$
the last bit of the output of $c$ (empty if that output is empty). Applying the
redirected action of $a$ with register $\rho$ (the moves and writes of $a$, no emission,
register updated to $a$'s emitted bit if any, and on a halt of $a$ either halting when
the updated register equals $b$ or entering the loop state otherwise) to the redirected
image of $c$ (the tapes and heads of $c$, empty output, and state determined by $c$'s
state and the last bit of its output in the same way) gives the redirected image of the
configuration obtained by applying $a$ to $c$.
\end{lemma}

\begin{lemma}[One-step correspondence for the halt-redirected machine]
\lean{Turing.FinTM.redirect_step}
\label{Turing.FinTM.redirect_step}
\leanok
\uses{Turing.Cfg, Turing.Cfg.inputSymbol, Turing.Cfg.workTapeSymbols, Turing.FinTM, Turing.FinTM.redirectAction, Turing.FinTM.redirectCfg, Turing.FinTM.redirectState, Turing.FinTM.redirectTM, Turing.FinTM.redirect_apply, Turing.FinTM.redirect_loop, Turing.MultiTapeTM.step, Turing.MultiTapeTM.step_of_halt}
\difficulty{5}
Let $M$ be a finite-state Turing machine over the Boolean alphabet and $b$ a bit, and
consider the halt-redirected machine of $M$ and $b$, which simulates $M$ silently while
recording the last bit $M$ has emitted, halts when $M$ halts with that bit equal to $b$,
and otherwise enters an idle loop state. For a configuration $c$ of $M$ on an input $x$,
let its redirected image carry the tapes and heads of $c$, empty output, and the state
$(q, \rho)$ when $c$ is in state $q$ and $\rho$ is the last bit of its output (empty if
none), or, when $c$ is halted, the halting state if $\rho = b$ and the loop state
otherwise. Then for every such $c$, halted or not, one step of the halt-redirected
machine from the redirected image of $c$ gives the redirected image of one step of $M$
from $c$.
\end{lemma}

\begin{lemma}[Run correspondence for the halt-redirected machine]
\lean{Turing.FinTM.redirect_run}
\label{Turing.FinTM.redirect_run}
\leanok
\uses{Turing.FinTM, Turing.FinTM.redirectCfg, Turing.FinTM.redirectTM, Turing.FinTM.redirect_step, Turing.MultiTapeTM.initCfg, Turing.MultiTapeTM.runFrom, Turing.MultiTapeTM.runFrom_comm_of_step}
\difficulty{2}
Let $M$ be a finite-state Turing machine over the Boolean alphabet and $b$ a bit, and
consider the halt-redirected machine of $M$ and $b$, which simulates $M$ silently while
recording the last bit $M$ has emitted, halts when $M$ halts with that bit equal to $b$,
and otherwise enters an idle loop state. For every input $x$ and every
$t \in \mathbb{N}$, the configuration of the halt-redirected machine after $t$ steps from
its initial configuration on $x$ is the redirected image of the configuration of $M$
after $t$ steps from its initial configuration on $x$: the same tapes and heads, empty
output, and state $(q, \rho)$ if $M$ is in state $q$ and the last bit it has emitted is
$\rho$ (empty if none), or, if $M$ has halted, halted when $\rho = b$ and in the loop
state otherwise.
\end{lemma}

\begin{theorem}[Halt-redirect: the halting clause]
\lean{Turing.FinTM.redirectTM_computes}
\label{Turing.FinTM.redirectTM_computes}
\leanok
\uses{Turing.FinTM, Turing.FinTM.ComputesInTime, Turing.FinTM.computesInTime_iff, Turing.FinTM.redirectCfg, Turing.FinTM.redirectState, Turing.FinTM.redirectTM, Turing.FinTM.redirect_run}
\difficulty{3}
Let $M$ be a finite-state Turing machine over the Boolean alphabet, $b$ a bit, $x$ and $w$
Boolean words, and $t \in \mathbb{N}$. Consider the halt-redirected machine of $M$ and
$b$, which simulates $M$ without emitting anything while recording the last bit $M$ has
emitted, halts when $M$ halts with that bit equal to $b$, and otherwise enters an idle
loop state. If $M$ on input $x$ halts within $t$ steps with output $w$, and $w$ is
nonempty with last bit $b$, then the halt-redirected machine on input $x$ halts within $t$
steps with empty output.
\end{theorem}

\begin{theorem}[Halt-redirect: the live clause]
\lean{Turing.FinTM.redirectTM_live}
\label{Turing.FinTM.redirectTM_live}
\leanok
\uses{Turing.Cfg.Halted, Turing.FinTM, Turing.FinTM.ComputesInTime, Turing.FinTM.ComputesInTime.output_unique, Turing.FinTM.computesInTime_iff, Turing.FinTM.redirectState, Turing.FinTM.redirectTM, Turing.FinTM.redirect_run, Turing.MultiTapeTM.initCfg, Turing.MultiTapeTM.runFrom}
\difficulty{4}
Let $M$ be a finite-state Turing machine over the Boolean alphabet, $b$ a bit, $x$ and $w$
Boolean words, and $t \in \mathbb{N}$. Consider the halt-redirected machine of $M$ and
$b$, which simulates $M$ without emitting anything while recording the last bit $M$ has
emitted, halts when $M$ halts with that bit equal to $b$, and otherwise enters an idle
loop state. If $M$ on input $x$ halts within $t$ steps with output $w$, and either $w$ is
empty or its last bit differs from $b$, then for every $u \in \mathbb{N}$ the
halt-redirected machine, run for $u$ steps from its initial configuration on $x$, is not
halted.
\end{theorem}

\begin{definition}[Decider padded with idle tapes]
\lean{Turing.FinTM.timedPadTM}
\label{Turing.FinTM.timedPadTM}
\leanok
\uses{Turing.FinTM, Turing.FinTM.leftAction, Turing.MultiTapeTM}
Let $D$ be a finite-state Turing machine over the Boolean alphabet with $k$ work tapes,
state set $Q$ and initial state $q_0$, and let $m \in \mathbb{N}$. The padded machine has
$k + m$ work tapes, state set $Q$ and initial state $q_0$; in state $q$, reading input
symbol $\sigma$ and work-tape symbols $w_0, \dots, w_{k+m-1}$, it performs $D$'s action
at $(q, \sigma, w_0, \dots, w_{k-1})$ on the input head and the first $k$ work tapes, with
the same emission and successor state, while writing nothing on the last $m$ work tapes
and keeping their heads still.
\end{definition}

\begin{definition}[Timed conditional controller]
\lean{Turing.FinTM.timedCondTM}
\label{Turing.FinTM.timedCondTM}
\leanok
\uses{Turing.FinTM, Turing.FinTM.branchTM, Turing.FinTM.controlAction, Turing.FinTM.leftAction, Turing.FinTM.rightAction, Turing.FinTM.timedPadTM, Turing.captureAction}
Let $D$, $M_1$, $M_2$ be finite-state Turing machines over the Boolean alphabet with
$k_D$, $k_1$, $k_2$ work tapes and state sets $Q_D$, $Q_1$, $Q_2$. The timed conditional
controller has $k_D + k_1 + k_2 + 1$ work tapes, split into a decider block of $k_D$
tapes, a branch block of $k_1 + k_2$ tapes and one capture tape; its states are those of
$Q_D$, two administrative states (back and read), a rewind-start and a rewind-scan state
for each bit $\beta$, and those of $Q_1 \sqcup Q_2$; and its initial state is $D$'s. In a
state of $D$ it performs $D$'s transition on the input and the decider block, writes $D$'s
emitted bit (if any) on the capture tape and moves that head right, emits nothing, leaves
the branch block idle, and enters the back state where $D$ would halt; the back state
moves the capture head one cell left and enters the read state; the read state reads the
bit $\beta$ under the capture head ($\mathrm{false}$ if blank) and enters rewind-start
for $\beta$; rewind-start moves the input head left and enters rewind-scan; rewind-scan
keeps moving the input head left while it reads an input symbol and, on reading the
blank boundary, moves right and enters the initial state of $M_1$ if $\beta$ is
$\mathrm{true}$ and of $M_2$ otherwise (all four administrative states writing and
emitting nothing and keeping the other heads still); and in a state of $M_i$ it performs
$M_i$'s transition, including its emission and halting, on the input and on $M_i$'s own
$k_i$ tapes within the branch block, leaving all other work tapes idle.
\end{definition}

\begin{definition}[Controller configuration in the branch phase]
\lean{Turing.FinTM.timedBranchCfg}
\label{Turing.FinTM.timedBranchCfg}
\leanok
\uses{Turing.Cfg, Turing.FinTM, Turing.FinTM.bufferTape, Turing.FinTM.leftCfg, Turing.FinTM.rightCfg, Turing.FinTM.timedCondTM}
Let $D$, $M_1$, $M_2$ be finite-state Turing machines over the Boolean alphabet with
$k_D$, $k_1$, $k_2$ work tapes, and consider the timed conditional controller for them,
whose work tapes form a decider block of $k_D$ tapes, a branch block of $k_1 + k_2$
tapes and one capture tape, and whose states include a copy of the states of $M_1$ and
$M_2$. Let $c$ be a configuration on an input $x$ of the two-branch union machine (which
runs $M_1$ on the first $k_1$ and $M_2$ on the last $k_2$ of its $k_1 + k_2$ work tapes,
according to whether its state belongs to $M_1$ or $M_2$), let $\tau$ and $h$ be
contents and head positions for $k_D$ tapes, and let $\beta$ be a bit. The branch-phase
configuration is the controller configuration on $x$ whose state is the copy of $c$'s
state (halted if $c$ is halted), whose input head and output are those of $c$, whose
decider block holds $\tau$ with heads $h$, whose branch block holds $c$'s work tapes and
heads, and whose capture tape holds the bit $\beta$ in cell $0$, blank elsewhere, with its
head at cell $0$.
\end{definition}

\begin{definition}[Controller configuration in the decider phase]
\lean{Turing.FinTM.timedControlCfg}
\label{Turing.FinTM.timedControlCfg}
\leanok
\uses{Turing.Cfg, Turing.FinTM, Turing.FinTM.leftCfg, Turing.FinTM.timedCondTM, Turing.captureCfg}
Let $D$, $M_1$, $M_2$ be finite-state Turing machines over the Boolean alphabet with
$k_D$, $k_1$, $k_2$ work tapes, and consider the timed conditional controller for them,
whose work tapes form a decider block of $k_D$ tapes, a branch block of $k_1 + k_2$
tapes and one capture tape, and whose states include a copy of the states of $D$ and an
administrative back state. For a configuration $c$ of $D$ on an input $x$, with output
$u$, the decider-phase configuration is the controller configuration on $x$ whose state
is the copy of $c$'s state, or the back state if $c$ is halted; whose input head is that
of $c$; whose decider block holds $c$'s work tapes and heads; whose branch block is blank
with every head at cell $0$; whose capture tape holds $u$ in cells $0, \dots, |u|-1$,
blank elsewhere, with its head at cell $|u|$; and whose output is empty.
\end{definition}

\begin{lemma}[Decider phase of the conditional controller]
\lean{Turing.FinTM.timed_capture}
\label{Turing.FinTM.timed_capture}
\leanok
\uses{Turing.Cfg, Turing.Cfg.Halted, Turing.FinTM, Turing.FinTM.leftCfg, Turing.FinTM.leftCfg_run, Turing.FinTM.timedCondTM, Turing.FinTM.timedControlCfg, Turing.FinTM.timedPadTM, Turing.MultiTapeTM.runFrom, Turing.capture_run}
\difficulty{4}
Let $D$, $M_1$, $M_2$ be finite-state Turing machines over the Boolean alphabet, and
consider the timed conditional controller for them, which first runs $D$ silently on a
decider block of tapes while copying $D$'s output onto a capture tape. Let $c$ be a
configuration of $D$ on an input $x$ and $t \in \mathbb{N}$ such that $D$ run from $c$ is
not halted at any time $s < t$. Then running the controller for $t$ steps from the
decider-phase configuration of $c$ (state of $c$, or the back state if $c$ is halted;
$c$'s input head, work tapes and heads on the decider block; a blank branch block; $c$'s
output stored from cell $0$ of the capture tape with that head just past it; empty
output) yields the decider-phase configuration of the configuration $D$ reaches after $t$
steps from $c$.
\end{lemma}

\begin{lemma}[Initial configuration of the conditional controller]
\lean{Turing.FinTM.timed_control_init}
\label{Turing.FinTM.timed_control_init}
\leanok
\uses{Turing.Cfg.init, Turing.FinTM, Turing.FinTM.leftCfg, Turing.FinTM.timedCondTM, Turing.FinTM.timedControlCfg, Turing.MultiTapeTM.initCfg, Turing.captureCfg}
\difficulty{4}
Let $D$, $M_1$, $M_2$ be finite-state Turing machines over the Boolean alphabet, and
consider the timed conditional controller for them, with a decider block, a branch
block and a capture tape as work tapes. For every input $x$, the initial configuration of
the controller on $x$ equals the decider-phase configuration of $D$'s initial
configuration on $x$: $D$'s initial state, the initial input-head position, every work
tape of all three blocks blank with its head at cell $0$, and empty output.
\end{lemma}

\begin{lemma}[Branch phase of the conditional controller]
\lean{Turing.FinTM.timed_branch_run}
\label{Turing.FinTM.timed_branch_run}
\leanok
\uses{Turing.Action.apply, Turing.Cfg, Turing.Cfg.inputSymbol, Turing.Cfg.workTapeSymbols, Turing.FinTM, Turing.FinTM.branchTM, Turing.FinTM.bufferTape, Turing.FinTM.leftAction, Turing.FinTM.leftCfg, Turing.FinTM.leftCfg_apply, Turing.FinTM.rightAction, Turing.FinTM.rightCfg, Turing.FinTM.rightCfg_apply, Turing.FinTM.timedBranchCfg, Turing.FinTM.timedCondTM, Turing.MultiTapeTM.runFrom, Turing.MultiTapeTM.runFrom_comm_of_step, Turing.MultiTapeTM.step}
\difficulty{5}
Let $D$, $M_1$, $M_2$ be finite-state Turing machines over the Boolean alphabet with
$k_D$, $k_1$, $k_2$ work tapes, consider the timed conditional controller for them, and
for a bit $\beta$ consider the two-branch union machine that runs $M_1$ or $M_2$ on
disjoint blocks of $k_1$ and $k_2$ work tapes and starts in the initial state of $M_1$ if
$\beta$ is $\mathrm{true}$ and of $M_2$ otherwise. Let $c$ be a configuration of the
union machine on an input $x$, $\tau$ and $h$ contents and head positions for $k_D$
tapes, and $t \in \mathbb{N}$. Then running the controller for $t$ steps from the
branch-phase configuration of $(c, \tau, h, \beta)$ (the state, input head, output and
work tapes of $c$ on the branch block, $\tau$ and $h$ on the decider block, and the bit
$\beta$ in cell $0$ of the capture tape with that head at cell $0$) yields the
branch-phase configuration of $(c', \tau, h, \beta)$, where $c'$ is the configuration the
union machine reaches after $t$ steps from $c$.
\end{lemma}

\begin{definition}[Controller configuration ready to rewind]
\lean{Turing.FinTM.timedReadyCfg}
\label{Turing.FinTM.timedReadyCfg}
\leanok
\uses{Turing.Cfg, Turing.FinTM, Turing.FinTM.branchTM, Turing.FinTM.timedBranchCfg, Turing.FinTM.timedCondTM, Turing.MultiTapeTM.initCfg}
Let $D$, $M_1$, $M_2$ be finite-state Turing machines over the Boolean alphabet, and
consider the timed conditional controller for them, whose work tapes form a decider block,
a branch block and a capture tape, and whose states include a rewind-start state for
each bit. For a configuration $c$ of $D$ on an input $x$ and a bit $\beta$, the
ready-to-rewind configuration is the controller configuration on $x$ in the rewind-start
state for $\beta$, with the input head of $c$, the decider block holding $c$'s work tapes
and heads, the branch block blank with every head at cell $0$, the capture tape holding
the bit $\beta$ in cell $0$, blank elsewhere, with its head at cell $0$, and empty output.
\end{definition}

\begin{lemma}[Reading the decider's one-bit verdict]
\lean{Turing.FinTM.timed_read}
\label{Turing.FinTM.timed_read}
\leanok
\uses{Turing.Action, Turing.Action.apply, Turing.Cfg, Turing.Cfg.workTapeSymbols, Turing.FinTM, Turing.FinTM.branchTM, Turing.FinTM.bufferTape, Turing.FinTM.controlAction, Turing.FinTM.controlAction_apply, Turing.FinTM.leftCfg, Turing.FinTM.rightCfg, Turing.FinTM.timedBranchCfg, Turing.FinTM.timedCondTM, Turing.FinTM.timedControlCfg, Turing.FinTM.timedReadyCfg, Turing.MultiTapeTM.initCfg, Turing.MultiTapeTM.runFrom, Turing.MultiTapeTM.step, Turing.captureCfg, Turing.moveInputPos_zero}
\difficulty{5}
Let $D$, $M_1$, $M_2$ be finite-state Turing machines over the Boolean alphabet, and
consider the timed conditional controller for them, with a decider block, a branch block
and a capture tape as work tapes. Let $c$ be a halted configuration of $D$ on an input
$x$ whose output is the one-bit word $\beta$. Then two steps of the controller from the
decider-phase configuration of $c$ (in the back state, with $c$'s input head and work
tapes on the decider block, a blank branch block, $\beta$ in cell $0$ of the capture tape
with that head at cell $1$, and empty output) reach the ready-to-rewind configuration of
$c$ and $\beta$, which differs only in being in the rewind-start state for $\beta$ with the
capture head at cell $0$.
\end{lemma}

\begin{lemma}[Reaching the selected branch within twice the decider budget]
\lean{Turing.FinTM.timed_start}
\label{Turing.FinTM.timed_start}
\leanok
\uses{Turing.Cfg.init, Turing.FinTM, Turing.FinTM.ComputesInTime, Turing.FinTM.ComputesInTime.output_unique, Turing.FinTM.branchTM, Turing.FinTM.computesInTime_iff, Turing.FinTM.timedBranchCfg, Turing.FinTM.timedCondTM, Turing.FinTM.timedControlCfg, Turing.FinTM.timedReadyCfg, Turing.FinTM.timed_capture, Turing.FinTM.timed_control_init, Turing.FinTM.timed_read, Turing.FinTM.timed_rewind, Turing.MultiTapeTM.initCfg, Turing.MultiTapeTM.runFrom, Turing.MultiTapeTM.runFrom_add, Turing.MultiTapeTM.timed_input_bound}
\difficulty{5}
Let $D$, $M_1$, $M_2$ be finite-state Turing machines over the Boolean alphabet with
$k_D$, $k_1$, $k_2$ work tapes, consider the timed conditional controller for them, let
$x$ be a Boolean word, $\beta$ a bit and $T \in \mathbb{N}$, and suppose $D$ on input $x$
halts within $T$ steps with output the one-bit word $\beta$. Then there exist
$a \le 2T + 5$ and contents $\tau$ and head positions $h$ for $k_D$ tapes such that after
$a$ steps from its initial configuration on $x$ the controller is in the branch-phase
configuration built from the initial configuration on $x$ of the two-branch union machine
selecting $\beta$: in the copy of the initial state of $M_1$ if $\beta$ is $\mathrm{true}$
and of $M_2$ otherwise, with the initial input-head position, a blank branch block with
every head at cell $0$, the decider block holding $\tau$ with heads $h$, the bit $\beta$ in
cell $0$ of the capture tape with that head at cell $0$, and empty output.
\end{lemma}

\begin{theorem}[Timed conditional composition of machines]
\lean{Turing.FinTM.computesFunInTime_cond}
\label{Turing.FinTM.computesFunInTime_cond}
\leanok
\uses{Turing.FinTM, Turing.FinTM.ComputesFunInTime, Turing.FinTM.ComputesInTime, Turing.FinTM.ComputesInTime.mono, Turing.FinTM.branchTM, Turing.FinTM.branchTM_computes, Turing.FinTM.computesInTime_iff, Turing.FinTM.leftCfg, Turing.FinTM.rightCfg, Turing.FinTM.timedBranchCfg, Turing.FinTM.timedCondTM, Turing.FinTM.timed_branch_run, Turing.FinTM.timed_start, Turing.MultiTapeTM.runFrom_add}
\difficulty{4}
Let $D$, $M_1$, $M_2$ be finite-state Turing machines over the Boolean alphabet, let $p$
be a Boolean predicate on Boolean words, $f_1$ and $f_2$ functions from Boolean words to
Boolean words, and $T_0, T_1, T_2 : \mathbb{N} \to \mathbb{N}$. Suppose that on every input
$x$, $D$ halts within $T_0(|x|)$ steps with output the one-bit word $p(x)$, $M_1$ halts
within $T_1(|x|)$ steps with output $f_1(x)$, and $M_2$ halts within $T_2(|x|)$ steps with
output $f_2(x)$. Then there exist a finite-state Turing machine $M$ over the Boolean
alphabet and a constant $c \in \mathbb{N}$ such that on every input $x$, $M$ halts within
$c \cdot (T_0(|x|) + \max(T_1(|x|), T_2(|x|)) + 1)$ steps with output $f_1(x)$ if $p(x)$
holds and $f_2(x)$ otherwise.
\end{theorem}
